\documentclass[12pt]{article}
\usepackage[utf8]{inputenc}
\usepackage[T1]{fontenc}
\usepackage{hyperref}
\title{Matu Daro Melanau Language Resources for K-12 Curriculum Development}
\author{Compiled by Claude Sonnet 5.5 (claude-sonnet-5-5) for Andr\'as Kornai}
\date{}
\begin{document}
\maketitle

\section{Introduction}
Daro-Matu Melanau (ISO 639-3 \texttt{dro}; Glottocode daro1239) is an Austronesian language of coastal Sarawak, Malaysia, closely related to the wider Melanau group. Wikipedia cites about 7,600 speakers (1981); Joshua Project gives a much larger figure (23,000) for the Daro-Matu people, so figures conflict. Little dedicated material was found.

\section{Primary Educational Resources}
No primers or textbooks were found.

\section{Linguistic and Reference Materials}
\begin{itemize}
\item Wikipedia on Daro-Matu: \url{https://en.wikipedia.org/wiki/Daro-Matu_language}
\item Wikipedia on the Melanau language group: \url{https://en.wikipedia.org/wiki/Melanau_language}
\item Bernama news report (2019) that a Melanau--Malay dictionary covering the Matu, Daro and Pulau Bruit areas was about 60\% complete under a UNESCO, Dewan Bahasa dan Pustaka and Ministry of Education project; I could not confirm it was published: \url{https://www.bernama.com/en/news.php?id=1774706}
\end{itemize}

\section{Religious and Community Resources}
\begin{itemize}
\item Joshua Project profile of the Daro-Matu in Malaysia: \url{https://joshuaproject.net/people_groups/11509}
\end{itemize}

\section{Digital and Audio Resources}
None found.

\section{Curriculum Development Context}
The language is not known to be taught in schools. The dictionary project reported by Bernama is the most relevant lead for curriculum work and should be checked directly with the organisations involved.

\end{document}